% Generated draft: requires mathematical and source review.
\documentclass[11pt]{article}
\usepackage[margin=25mm]{geometry}
\usepackage{fontspec}
\setmainfont{Latin Modern Roman}
\usepackage{amsmath,amssymb,mathtools,unicode-math}
\setmathfont{Latin Modern Math}
\usepackage{enumitem,needspace,xurl,hyperref}
\setlength{\parindent}{0pt}
\setlength{\parskip}{5pt}
\setlength{\emergencystretch}{3em}
\newcommand{\N}{\mathbb{N}}
\newcommand{\Z}{\mathbb{Z}}
\newcommand{\Q}{\mathbb{Q}}
\newcommand{\R}{\mathbb{R}}
\newcommand{\C}{\mathbb{C}}
\newcommand{\F}{\mathbb{F}}
\newcommand{\A}{\mathbb{A}}
\newcommand{\PP}{\mathbb{P}}
\newcommand{\E}{\mathbb{E}}
\newcommand{\eps}{\varepsilon}
\newcommand{\dd}{\,\mathrm{d}}
\begin{document}
\section*{UnsolvedMath ID: 30001252}
\textbf{Source problem code:} catalogue entry
\subsection{Definitions and mathematical statement}
Let \(h>0\), and let \((a_n)_{n\geq1}\) be the prescribed sequence of nonnegative masses in the supplied record. Retaining its notation, set
\[
A(h,(a_n))=\left\{\mu\in M_+(\R_+):\ \mu([nh,(n+1)h])=a_n\ \text{for every integer }n\geq1\right\}.
\]
Here \(M_+(\R_+)\) denotes the source's class of nonnegative measures on \(\R_+\); the supplied record does not specify the measurable structure, finiteness or local finiteness requirements, or the convention for \(\R_+\).

For a nonempty set \(B\) on which a distance \(\rho\) is defined, its diameter is
\[
\operatorname{diam}_{\rho}(B)=\sup\{\rho(\mu,\nu):\mu,\nu\in B\},
\]
allowing the supremum to be \(+\infty\). The requested problem is to determine
\[
\operatorname{diam}_{\rho_F}\bigl(A(h,(a_n))\bigr)
\quad\text{and}\quad
\operatorname{diam}_{\rho_W}\bigl(A(h,(a_n))\bigr),
\]
and describe their dependence on \(h\), with the prescribed sequence held fixed. The record calls \(\rho_F\) the flat metric and \(\rho_W\) a Wasserstein-type (also described as Wasserstein-like) metric, but supplies neither definition. Consequently, these expressions retain the source's symbols and do not specify unique distances without further information.

The intervals are closed exactly as supplied: adjacent intervals share endpoints, so an atom at a shared endpoint contributes to both corresponding constraints. The indexing begins at \(n=1\); no interval-mass condition is imposed on \(\R_+\cap[0,h)\). No summability, total-mass, moment, or nonemptiness hypothesis is supplied. These omissions must be resolved, together with the metric definitions and their domains, before the diameter questions have a fully specified interpretation; no convention for an empty-set diameter is imposed here.
\subsection{Short English statement}
Given a positive interval length \(h\) and prescribed nonnegative masses \(a_n\), consider all measures assigning mass \(a_n\) to each closed interval \([nh,(n+1)h]\), for every \(n\geq1\). Determine the diameter of this family in the source's flat and Wasserstein-type metrics and how each diameter depends on \(h\). The metric definitions, measure class, and necessary admissibility hypotheses are missing from the supplied record, so this remains a qualified statement requiring mathematical review.
\subsection{Sources}
\begin{itemize}
\item Local UnsolvedMath registry, supplied complete source record, permanent ID \(30001252\), source problem code \texttt{catalogue entry}, title: \emph{Metric Diameters of Measures with Prescribed Interval Masses}. The registry assigns proposed year 2009 and supplies no named proposer.
\item Supplied report citation: \emph{Mathematical Biology} (2009), \emph{Oberwolfach Reports}, pp.\ 1303--1374, DOI \texttt{10.4171/owr/2009/24}. Source URL and supplied literature-source link, labeled ``original Oberwolfach report'': \url{https://doi.org/10.4171/owr/2009/24}. The background describes Oberwolfach Reports as documenting mathematical workshops and their open-problem discussions. No problem-specific page or metric definitions are supplied.
\item Supplied literature assessment, checked on \texttt{2026-08-21}: the registry labels the problem open and states that the report gives upper bounds for several metrics while leaving the exact diameter and its dependence on \(h\) undetermined. This is supplied dated curation derived from the OWR record, not an independently verified literature assessment; time-sensitive status and cited claims require independent verification.
\end{itemize}
\end{document}
